\documentclass[12pt]{article}
\usepackage[utf8]{inputenc}
\usepackage[T1]{fontenc}
\usepackage{hyperref}
\title{Tetun Language Resources for K-12 Curriculum Development}
\author{Compiled by Claude Sonnet 5 (claude-sonnet-5) for Andr\'as Kornai}
\date{}
\begin{document}
\maketitle

\section{Introduction}
Tetun (Tetum, ISO 639-3 tet) is a co-official language of Timor-Leste, alongside Portuguese. The search found dictionaries, a learner course and policy discussion of language in education. It did not find a current Ministry of Education Tetun textbook series, so that would need a separate search.

\section{Primary Educational Resources}
\begin{itemize}
\item \emph{Tetum, A Language For Everyone} (\emph{Tetun, lian ida ba ema hotu-hotu}), a course book for English speakers with dialogues, exercises, a phrasebook and a two-way dictionary (retailer listing): \url{https://universidades.umlivro.com.br/tetum--a-language-for-everyone--tetun--lian-ida-ba-ema-hotu-hotu--1630269/p}
\item ANU Tetum course listing (TETM6002), an intermediate-level university course: \url{https://programsandcourses.anu.edu.au/2019/course/TETM6002/First%20Semester/3443}
\end{itemize}

\section{Linguistic and Reference Materials}
\begin{itemize}
\item Tetun Dili dictionary on Webonary: \url{https://www.webonary.org/tetun-dili/overview/introduction/}
\item Cliff Morris, \emph{A Traveller's Dictionary in Tetun-English and English-Tetun}, based on the south-coast Tetun-Los dialect, about 1,100 words: \url{https://levlafayette.com/node/44}
\item The same dictionary, GNU Tetum project copy: \url{https://www.gnu.org/software/tetum/contributors/cliffMorris-xhtml/pr04.html}
\end{itemize}

\section{Religious and Community Resources}
No religious or community-organisation resources were found in the search.

\section{Digital and Audio Resources}
The Webonary dictionary is the only digital resource found. No audio or video resources turned up.

\section{Curriculum Development Context}
\begin{itemize}
\item East Timor Women Australia, \emph{Languages in Education in East Timor}: \url{https://www.etwa.org.au/wp-content/uploads/2013/11/languages_education.pdf}
\item UNHCR Refworld, \emph{Timor-Leste: When do mother tongues divide?}: \url{https://webarchive.archive.unhcr.org/20230518164944oe_/https://www.refworld.org/docid/4f7428972.html}
\end{itemize}
These sources discuss Tetun and Portuguese as co-official languages and the debate over mother-tongue-based multilingual education. They report that Tetum is a first language for only part of the population and a second or third language for most.

Dialect variation (Tetun Dili versus Tetun-Los) matters for choosing a standard.
\end{document}